% !TEX root = userguide.tex

\def\headdate{20th February 2018}

\section{Diffuse-interface model with particles}
\label{sec:dim_with_particles}

(These examples have been developed by Nick Jaensson of the TU/e)

\subsection{Particle-covered droplet sheared between two plates}

Example \texttt{dimpart1.f90} describes a Newtonian droplet in a Newtonian
matrix which is covered by rigid particles. The fluid-fluid interface is 
described by the Cahn-Hilliard diffuse-interface model (see Section 
\ref{sec:examples_dim}). The particle-fluid interface is described by a
sharp-interface model using boundary-fitted meshes and an ALE formulation 
(see Section \ref{sec:ale}). To
include the effect of surface tension, the stress-form of the Cahn-Hilliard
theory is used \cite{Jaensson2015}. The (static) contact angle between the
fluid-fluid interface and the particle boundaries can be adapted. The example 
makes extensive use of adaptive meshing using refinement points, as explained
in Section \ref{sec:examples_dim}. The example can be obtained by typing at the
command line:
\begin{quote}
   \texttt{getexample dim\_with\_particles}
\end{quote}

In Figure~\ref{fig:dimpart} the result of a sheared particle-covered droplet
(initially circular) after a number of time steps is shown. Note, how the
particles are located ``on the outside'' of the droplet due to the imposed 
contact angle.
\begin{figure}[htbp!]
   \begin{center}
   \includegraphics[scale=0.3]{dimpart}
   \end{center}
   \caption{Particle-covered droplet sheared between two plates.
     Results obtained using example \texttt{dimpart1.f90}.}
    \label{fig:dimpart}
\end{figure}
